\documentclass[11pt,a4paper]{article}

% --- Packages ---
\usepackage[utf8]{inputenc}
\usepackage[T1]{fontenc}
\usepackage{lmodern}
\usepackage[margin=2.5cm]{geometry}
\usepackage{amsmath,amssymb}
\usepackage{graphicx}
\usepackage{hyperref}
\usepackage{booktabs}
\usepackage{xcolor}
\usepackage{listings}
\usepackage{algorithm}
\usepackage{algpseudocode}
\usepackage{natbib}
\usepackage{microtype}
\usepackage{enumitem}
\usepackage{titlesec}
\usepackage{abstract}

\hypersetup{
  colorlinks=true,
  linkcolor=blue!60!black,
  citecolor=green!50!black,
  urlcolor=blue!60!black
}

\lstset{
  basicstyle=\small\ttfamily,
  breaklines=true,
  frame=single,
  rulecolor=\color{gray!40},
  backgroundcolor=\color{gray!5}
}

% --- Title ---
\title{%
  \textbf{Directive Grounding as a Substitute for Model Scale:}\\
  \large How Structured Skill Definitions Enable Quantized LLMs\\
  to Perform Specialist One-Shot Tasks
}

\author{
  The Agency Framework\\
  \texttt{https://github.com/flegare/the\_agency}
}

\date{March 2026}

% ===========================================================================
\begin{document}
\maketitle

% --- Abstract ---
\begin{abstract}
Large language models (LLMs) are commonly scaled up—in both parameter count
and numerical precision—to handle complex, open-ended reasoning tasks. We
challenge the assumption that \emph{model scale} is the primary driver of
task-completion quality in structured, domain-specific agentic workflows.
Instead, we propose the \textbf{Directive Grounding Hypothesis}: that the
behavioral intelligence required to complete a well-scoped software
engineering task can be largely externalized into a structured persona
definition—a \emph{Skill}—thereby reducing the residual cognitive burden
placed on the underlying model to a level achievable by aggressively
quantized, locally-executable LLMs. We formalize a model of intelligence
distribution across the skill definition and model weights, derive an
operational criterion for skill directiveness, and demonstrate through the
design of \textit{The Agency} framework how this principle yields reliable
one-shot task completion using 3--8~billion parameter quantized models where
ungrounded prompting fails. Our findings suggest a practical path toward
democratizing capable AI-assisted software development without dependence on
proprietary cloud APIs.
\end{abstract}

\tableofcontents
\newpage

% ===========================================================================
\section{Introduction}

The dominant paradigm in large language model deployment assumes a direct
relationship between model capability and model scale. Practitioners
routinely reach for larger, higher-precision models when facing complex tasks:
tasks that smaller models fail at in isolation. This creates a structural
dependency on cloud-hosted APIs and GPU clusters far beyond the reach of
individual developers or resource-constrained organizations.

At the same time, quantization research has demonstrated that models can be
compressed to 4-bit or even 3-bit precision with surprisingly small losses in
general benchmark performance~\citep{frantar2022gptq, dettmers2022llmint8}.
Yet practitioners report that quantized models struggle with the kinds of
multi-step, open-ended reasoning tasks that characterize real software
engineering work—not because their weights are insufficient in isolation, but
because the \emph{implicit problem specification} they receive leaves too much
ambiguity for a compressed reasoning engine to resolve reliably.

We argue that this failure mode is a property of \emph{task formulation}, not
model capacity. A quantized 7B model presented with the instruction
\textit{``write a secure authentication module''} must simultaneously:

\begin{enumerate}[noitemsep]
  \item Interpret the domain (web auth, CLI auth, API keys, OAuth\ldots)
  \item Select a security posture (OWASP, zero-trust, minimal attack surface\ldots)
  \item Infer the project's technology stack
  \item Choose a code structure and naming convention
  \item Decide on test coverage requirements
  \item Determine deployment constraints
\end{enumerate}

Each of these sub-decisions consumes representational capacity from the model.
Crucially, none of them are \emph{reasoning} tasks—they are
\emph{disambiguation} tasks. If the task specification were precise enough to
eliminate all ambiguity, the model could direct its full capacity toward code
generation itself.

This is the central insight of this paper: \textbf{a sufficiently directive
skill definition can absorb the disambiguation burden, leaving only the
residual creative and syntactic synthesis task to the model}. The required
model capacity for the residual task is substantially lower than for the
combined disambiguation-plus-synthesis task—potentially low enough to be
handled by quantized local models.

The remainder of this paper is structured as follows. Section~\ref{sec:background}
reviews relevant work on quantization, prompt engineering, and agent frameworks.
Section~\ref{sec:hypothesis} formalizes the Directive Grounding Hypothesis.
Section~\ref{sec:intelligence-dist} develops a model of intelligence
distribution. Section~\ref{sec:directiveness} introduces a practical criterion
for skill directiveness. Section~\ref{sec:agency} describes how The Agency
framework instantiates these principles. Section~\ref{sec:implications}
discusses implications for deployment. Section~\ref{sec:future} outlines an
experimental validation program. Section~\ref{sec:conclusion} concludes.

% ===========================================================================
\section{Background}
\label{sec:background}

\subsection{Model Quantization}

Quantization reduces the numerical precision of model weights from 32-bit or
16-bit floating point to lower-bit integer representations. The seminal work
of \citet{dettmers2022llmint8} demonstrated that 8-bit quantization of large
models could be achieved with minimal accuracy loss using mixed-precision
decomposition. \citet{frantar2022gptq} extended this to 4-bit precision using
approximate second-order weight updates, enabling 70B parameter models to run
on consumer hardware.

Subsequent work such as GGUF/llama.cpp~\citep{gerganov2023llamacpp} and
Ollama~\citep{ollama2023} has made quantized local inference accessible to
practitioners without specialized knowledge, typically at 3--5 bits per weight.

The standard critique of quantized models is degraded performance on
``reasoning-heavy'' benchmarks (e.g., GSM8K, HumanEval). However, we note
that these benchmarks measure \emph{ungrounded} reasoning: the model must
derive all context from first principles. This is a maximally unfavorable
evaluation setting for quantized models.

\subsection{Prompt Engineering and Chain-of-Thought}

Research on chain-of-thought prompting~\citep{wei2022chain} demonstrated that
decomposing complex reasoning tasks into explicit intermediate steps
dramatically improves the performance of smaller models. This is an early
instance of the principle we formalize: externalizing cognitive structure
reduces the load on model weights. Least-to-most prompting~\citep{zhou2022least}
and program-of-thought~\citep{chen2022program} extend this principle further.

Zero-shot chain-of-thought~\citep{kojima2022large} showed that even a simple
directive (``Let's think step by step.'') could unlock substantial latent
capability in mid-sized models. If a four-word directive achieves this,
structured multi-page persona definitions should achieve substantially more.

\subsection{Agent Frameworks and Role-Playing}

ReAct~\citep{yao2023react} demonstrated that interleaving reasoning and acting
in a structured loop improves LLM agent performance. Generative
Agents~\citep{park2023generative} showed that rich contextual scaffolding about
an agent's identity, memory, and social role produces coherent long-horizon
behavior from GPT-4-class models.

Crucially, \citet{chen2023agentverse} showed that \emph{role specialization}
in multi-agent systems—where each agent has a narrow, well-defined
responsibility—produces better outcomes than generalist agents, even when the
underlying models are identical. The performance gain derives from the role
definition, not the model.

Toolformer~\citep{schick2023toolformer} demonstrated that teaching a model
when and how to invoke external tools—a form of procedural grounding—
dramatically expands the effective task frontier of a 6.7B parameter model.
We interpret this as another instance of intelligence distribution: the tool
usage protocol absorbs the decision-making load, leaving the model to focus
on content generation.

\subsection{The Lost-in-the-Middle Problem}

\citet{liu2023lost} showed that LLMs suffer from significant positional bias
in long contexts: information in the middle of a long context window is
systematically underweighted relative to information at the beginning or end.
This finding has direct implications for skill design—the most critical
constraints must be placed at the beginning or end of the skill definition
to ensure they are reliably attended to by quantized models with smaller
effective context windows.

% ===========================================================================
\section{The Directive Grounding Hypothesis}
\label{sec:hypothesis}

\begin{description}
  \item[Definition (Grounding).] A \emph{grounding} $G$ for a task $T$ is a
  set of declarative and procedural statements that reduce the effective
  search space $\mathcal{S}(T)$ over valid outputs. More directive groundings
  have smaller residual search spaces.

  \item[Definition (Skill).] A \emph{Skill} $\mathcal{K}$ is a structured
  text document comprising: (i) a role identity, (ii) a set of core
  responsibilities, (iii) explicit workflow integration rules, and (iv) named
  constraints and heuristics applicable to the role's domain.

  \item[Definition (Residual Task).] Given a task $T$ and a skill
  $\mathcal{K}$, the \emph{residual task} $T | \mathcal{K}$ is the task
  remaining after applying the grounding provided by $\mathcal{K}$:
  \[
    T | \mathcal{K} = T \setminus \text{decisions}(\mathcal{K})
  \]
  where $\text{decisions}(\mathcal{K})$ denotes the set of sub-decisions
  resolved by the skill definition.
\end{description}

\vspace{0.5em}
\noindent\textbf{Directive Grounding Hypothesis.} \textit{There exists a skill
$\mathcal{K}^*$ for any bounded software engineering task $T$ such that the
residual task $T | \mathcal{K}^*$ lies within the reliable capability envelope
of an aggressively quantized local LLM $M_q$, even though $T$ itself does not:}
\[
  P(M_q \text{ succeeds on } T) \ll P(M_q \text{ succeeds on } T | \mathcal{K}^*)
\]

This hypothesis makes a strong claim: the bottleneck for quantized model
failure on complex tasks is not \emph{weight capacity} per se, but the
\emph{unresolved ambiguity} in the task specification. Resolving that
ambiguity externally, through a directive skill definition, unlocks latent
model capacity that was previously consumed by disambiguation.

\vspace{0.5em}
\noindent\textbf{Corollary (Quantization-Skill Substitutability).}
\textit{Within bounded domains, increasing skill directiveness is a
partial substitute for increasing model scale or precision.} That is, for
a target task-completion quality $Q^*$:
\[
  (M_{\text{large}}, \mathcal{K}_{\text{weak}}) \approx_Q
  (M_{\text{small}}, \mathcal{K}_{\text{directive}})
\]

This corollary is the operational claim of this paper: intelligence can be
traded across the model--skill boundary.

% ===========================================================================
\section{A Model of Intelligence Distribution}
\label{sec:intelligence-dist}

We model the total intelligence required to complete a task $T$ as a scalar
$I(T)$. This total is distributed across two sources:

\begin{enumerate}
  \item \textbf{Parametric intelligence} $I_\theta$: derived from the model's
  trained weights. This is the component reduced by quantization.
  \item \textbf{Contextual intelligence} $I_c$: derived from the information
  provided in the context window at inference time—including the skill
  definition, project state, and task description.
\end{enumerate}

For a task to succeed, we require:
\[
  I_\theta(M) + I_c(\mathcal{K}, \text{context}) \geq I(T)
\]

For a full-precision large model, $I_\theta$ is high, so the bar on $I_c$ is
low—a vague task description suffices. For a quantized small model,
$I_\theta$ is reduced, so $I_c$ must compensate. A directive skill definition
systematically maximizes $I_c$.

\subsection{Components of Contextual Intelligence}

A well-designed skill definition contributes to $I_c$ through four channels:

\begin{description}
  \item[Role Grounding.] Establishing a specific identity (e.g.,
  \textit{application\_security\_engineer}) primes the model's attention toward
  a narrow region of its parametric knowledge, reducing activation of
  irrelevant or conflicting patterns from other domains.

  \item[Responsibility Scoping.] An explicit numbered list of core
  responsibilities converts a complex open task into a constrained sequence
  of subtasks. Each subtask has a substantially smaller search space than
  the combined task.

  \item[Workflow Proceduralization.] Explicit handoff rules (``after
  completing X, pass to Y skill'') externalize the planning component of
  task execution. Planning is expensive for quantized models; proceduralization
  eliminates it.

  \item[Heuristic Constraints.] Named behavioral rules (e.g., ``Cyclic Error
  Watchdog: if stuck in an error loop for more than 3 iterations, stop and
  report'') encode domain expertise that would otherwise require the model to
  derive from first principles under uncertainty.
\end{description}

\subsection{The Disambiguation Ceiling}

Figure~\ref{fig:dist} illustrates the key relationship. As a task grows in
complexity (x-axis), the disambiguation ceiling for quantized models is
reached quickly—beyond a threshold, ungrounded models fail reliably. A
directive skill shifts the effective task complexity downward by absorbing
disambiguation load, moving the task back within the quantized model's
reliable operating region.

\begin{figure}[h]
\centering
\begin{minipage}{0.85\textwidth}
\lstset{frame=none,backgroundcolor=\color{white}}
\begin{lstlisting}[language={},basicstyle=\small\ttfamily]
Task Success Rate
    ^
1.0 |......................
    |              ......        Full-precision large model
    |          ....
    |       ...
0.5 |    ...               ....         Quantized + Directive Skill
    |   .               ...
    |  .             ...
    | .            ..
    |.           ..              Quantized, no skill grounding
0.0 +............._____________._____________> Task Complexity
         low         medium          high
\end{lstlisting}
\end{minipage}
\caption{Schematic relationship between task complexity and success rate for
three configurations: full-precision large model (top), quantized model with
directive skill grounding (middle), and quantized model with minimal grounding
(bottom). Directive grounding shifts the quantized model's capability curve
rightward, toward higher task complexity.}
\label{fig:dist}
\end{figure}

% ===========================================================================
\section{Measuring Skill Directiveness}
\label{sec:directiveness}

We propose a practical \emph{Directiveness Score} $D(\mathcal{K})$ for a
skill definition, computed as a weighted sum of structural properties:

\[
  D(\mathcal{K}) = w_1 \cdot R + w_2 \cdot P + w_3 \cdot H + w_4 \cdot B
\]

where:
\begin{itemize}[noitemsep]
  \item $R$ = number of explicitly numbered responsibilities
  \item $P$ = number of explicit procedural handoff rules
  \item $H$ = number of named heuristic constraints
  \item $B$ = 1 if the skill includes a ``Discovery Protocol'' (must read
    before writing), 0 otherwise
  \item $w_1, w_2, w_3, w_4$ are empirically determined weights
\end{itemize}

A skill with $D(\mathcal{K}) = 0$ is a purely identity-based prompt (e.g.,
``You are a software engineer.''). A skill with high $D(\mathcal{K})$
provides an exhaustive procedural specification. We hypothesize a monotonic
relationship between $D(\mathcal{K})$ and $P(M_q \text{ succeeds on }
T | \mathcal{K})$ for fixed $M_q$ and $T$.

\subsection{Anti-Patterns in Low-Directive Skills}

The following skill structures are associated with low directiveness and
high failure rates in quantized models:

\begin{description}
  \item[Identity-only prompts.] ``You are a security expert.'' Provides role
    grounding but no responsibility scoping or procedural structure.
  \item[Capability lists.] ``You can do X, Y, and Z.'' Describes capability
    without constraining behavior or sequencing.
  \item[Outcome statements.] ``Produce high-quality, secure code.'' Specifies
    a desired output quality without providing a path to achieve it.
\end{description}

\subsection{Patterns in High-Directive Skills}

Conversely, high-directive skills exhibit:

\begin{description}
  \item[Ordered responsibility lists.] Each responsibility is numbered and
    action-oriented (verb-first: ``Perform X'', ``Validate Y'', ``Generate Z'').
  \item[Explicit pre-conditions.] ``Before generating any output, you MUST
    read [file]. If [file] does not exist, halt and report its absence.''
  \item[Named watchdogs.] ``Cyclic Error Watchdog: if the same error recurs
    after 3 fix attempts, stop and explain the contradiction.''
  \item[Handoff protocols.] ``Upon completion, produce a structured handoff
    document for the [next\_skill] detailing [specific outputs].''
  \item[Scope delimiters.] Explicit statements of what the skill does
    \emph{not} do, preventing scope creep that consumes model context.
\end{description}

% ===========================================================================
\section{The Agency Framework as an Implementation}
\label{sec:agency}

The Agency\footnote{\url{https://github.com/flegare/the_agency}} is an
open-source framework that implements these principles through 32 specialized
skill definitions organized as an enterprise software development organization.
Each skill is defined as a structured Markdown document (\texttt{SKILL.md})
injected into the context window of whatever AI assistant or LLM the
practitioner uses.

\subsection{Skill Architecture}

Each \texttt{SKILL.md} follows a canonical structure:

\begin{lstlisting}[language={},caption={Canonical SKILL.md structure}]
---
name: Application Security Engineer
description: Code hardening and threat modeling
---

# Application Security Engineer

[Role identity and organizational context]

## Your Core Responsibilities
1. [Ordered, action-oriented responsibility]
2. [...]
N. [...]

## Workflow Integration
[Explicit handoff rules: who triggers this skill,
 what it produces, who receives its outputs]
\end{lstlisting}

\subsection{The Discovery Protocol Pattern}

Several high-stakes skills in The Agency implement a \emph{Discovery Protocol}
that explicitly forbids output generation before file reading:

\begin{lstlisting}[language={},caption={Discovery Protocol excerpt from historian/SKILL.md}]
DISCOVERY PROTOCOL (Mandatory):
Before writing ANY output, you MUST:
1. Read docs/cmo_state.md (if it exists). Mark as [MISSING] if absent.
2. Read docs/architecture/system_architecture.md. Mark as [MISSING] if absent.
3. List every file you have read at the top of your response.

NEVER infer, assume, or hallucinate project state.
If a file is missing, document its absence explicitly.
\end{lstlisting}

This protocol is a direct implementation of the heuristic constraint
pattern described in Section~\ref{sec:directiveness}. It eliminates the
hallucination failure mode—where quantized models confidently generate
plausible but incorrect project state—by making file reading a pre-condition
for any output.

\subsection{The Project Manager as a Router}

The \texttt{project\_manager} skill implements the Responsibility Scoping
pattern at a meta-level: its sole responsibility is task decomposition and
routing to appropriate specialist skills. This separates the \emph{planning}
function (which is disambiguation-heavy) from the \emph{execution} function
(which is synthesis-heavy), allowing different quantization levels to be
applied to each phase.

\subsection{Workflow Integration Graph}

The Agency's skills form a directed acyclic graph (DAG) of handoffs:

\[
  \text{project\_manager} \to \text{cmo\_analyst} \to \text{solution\_architect}
  \to \text{ciso} \to \text{backend/frontend} \to \text{chief\_test\_officer}
  \to \text{ssdlc\_manager} \to \text{historian}
\]

Each edge in the DAG represents an explicit handoff protocol defined in the
sending skill's Workflow Integration section. A quantized model executing any
node in this graph need not plan the overall workflow—the plan is encoded in
the skill definition. This is a direct implementation of the workflow
proceduralization component of $I_c$.

\subsection{Skill Directiveness Scores in The Agency}

Table~\ref{tab:scores} provides indicative directiveness scores for a
representative sample of Agency skills, using equal weights
($w_1 = w_2 = w_3 = w_4 = 1$):

\begin{table}[h]
\centering
\caption{Indicative Directiveness Scores for selected Agency skills.
$R$: responsibilities, $P$: handoff rules, $H$: named heuristics,
$B$: discovery protocol present.}
\label{tab:scores}
\begin{tabular}{lccccr}
\toprule
Skill & $R$ & $P$ & $H$ & $B$ & $D(\mathcal{K})$ \\
\midrule
\texttt{historian}                    & 8  & 4 & 3 & 1 & 16 \\
\texttt{cmo\_analyst}                 & 12 & 5 & 4 & 1 & 22 \\
\texttt{application\_security\_engineer} & 7 & 3 & 3 & 0 & 13 \\
\texttt{chief\_test\_officer}         & 6  & 4 & 2 & 0 & 12 \\
\texttt{frontend\_developer}          & 5  & 2 & 1 & 0 &  8 \\
\texttt{coder}                        & 4  & 2 & 2 & 0 &  8 \\
\texttt{solution\_architect}          & 5  & 3 & 0 & 0 &  8 \\
\texttt{project\_manager}             & 4  & 6 & 0 & 0 & 10 \\
\midrule
Mean & 6.4 & 3.6 & 1.9 & 0.25 & 12.1 \\
\bottomrule
\end{tabular}
\end{table}

The \texttt{cmo\_analyst} skill—which maintains the master project state
document and is the highest-stakes single point of ground truth—has the
highest directiveness score. This is consistent with the hypothesis: tasks
requiring the most ground-truth fidelity receive the most directive
specifications.

% ===========================================================================
\section{Implications for Practical Deployment}
\label{sec:implications}

\subsection{Democratizing Agentic Software Development}

If the Directive Grounding Hypothesis holds, it has significant practical
implications. Organizations can replace API dependencies on GPT-4-class or
Claude-class models for structured coding workflows with locally-executable
3--8B quantized models (e.g., \texttt{llama3.2:3b}, \texttt{qwen2.5-coder:7b},
\texttt{codellama:7b} via Ollama), provided they invest in directive skill
design rather than model scale.

The economics are substantial: a locally-hosted quantized model costs the
electricity of one workstation GPU, versus API costs that can reach hundreds
of dollars per month for intensive agentic workflows.

\subsection{The Skill-as-Compression Trade-off}

We note a dual relationship: a directive skill definition can be thought of
as a \emph{lossless compression} of domain expertise into structured text.
The model at inference time \emph{decompresses} this expertise by applying
it to the specific task. The compression ratio (skill definition length to
effective intelligence contributed) is a key metric for skill design quality.

A well-designed skill should be \emph{maximally dense}—every sentence should
resolve a decision that the model would otherwise have to make
independently—while remaining \emph{concise enough} to fit within the
effective context window of a quantized model without triggering the
lost-in-the-middle degradation documented by \citet{liu2023lost}.

\subsection{Skill Versioning and Model-Specific Tuning}

Different quantized models have different ``disambiguation tolerances''—the
amount of residual ambiguity they can handle before failure rates become
unacceptable. A practical deployment strategy involves:

\begin{enumerate}
  \item Profiling target model disambiguation tolerance on a reference task suite.
  \item Measuring the directiveness deficit between the model's tolerance and
    the task's inherent ambiguity.
  \item Augmenting skill definitions to cover the deficit, without exceeding
    the context budget.
\end{enumerate}

This creates a natural versioning system for skills: \texttt{SKILL.md} for
large models, \texttt{SKILL.compact.md} for 7B models, and
\texttt{SKILL.minimal.md} for 3B models—each with increasing directiveness
to compensate for decreasing parametric capacity.

\subsection{Security Implications}

A corollary of high directiveness is reduced jailbreak surface. A model
operating under an exhaustive set of procedural constraints has fewer
degrees of freedom through which an adversarial input can redirect behavior.
This aligns with the prompt injection defense motivation of The Agency's
Basalt Shield integration, which validates all skill definitions against
prompt injection attempts before merging.

% ===========================================================================
\section{Experimental Validation Framework}
\label{sec:future}

The Directive Grounding Hypothesis is an empirical claim that requires
systematic validation. We outline a proposed experimental program:

\subsection{Experiment 1: Directiveness vs. Success Rate}

\textbf{Design.} For a fixed set of software engineering tasks $\{T_i\}$
and a fixed quantized model $M_q$ (e.g., \texttt{llama3.2:3b} at 4-bit),
produce skill definitions at varying directiveness levels
$D_1 < D_2 < \ldots < D_k$ by progressively adding responsibilities,
handoff rules, and named constraints.

\textbf{Metric.} One-shot task completion rate, assessed by automated test
suite execution and human review of output artifacts.

\textbf{Hypothesis.} Task completion rate increases monotonically with
directiveness up to a saturation point, after which additional directives
provide diminishing returns.

\subsection{Experiment 2: Quantization-Skill Substitutability}

\textbf{Design.} Fix a target task-completion quality $Q^*$. Identify the
minimum model size required to achieve $Q^*$ with a minimal skill definition
($D \approx 0$). Then, progressively increase skill directiveness and measure
the reduction in minimum model size required to maintain $Q^*$.

\textbf{Metric.} Minimum viable parameter count as a function of $D(\mathcal{K})$.

\textbf{Hypothesis.} A well-designed directive skill can substitute for
approximately one order of magnitude of parameter count (e.g., enabling a
7B model to match a 70B model on the target task).

\subsection{Experiment 3: Discovery Protocol Ablation}

\textbf{Design.} Compare task completion rates for skills with and without
the Discovery Protocol requirement, specifically targeting tasks where
hallucinated project state would cause downstream failures.

\textbf{Hypothesis.} The Discovery Protocol produces a statistically
significant reduction in hallucination-caused failures, with the largest
effect size in quantized models.

\subsection{Benchmark Tasks}

Candidate tasks for validation, spanning the Agency skill graph:

\begin{itemize}[noitemsep]
  \item Implement a JWT authentication middleware from an architecture spec
  \item Generate a full test suite for a given API contract
  \item Produce a threat model for a described microservice topology
  \item Migrate a schema from PostgreSQL to a graph database model
  \item Write a CI/CD pipeline definition from infrastructure constraints
\end{itemize}

Each task has an objectively verifiable success criterion (tests pass,
schema validates, pipeline executes) making automated evaluation tractable.

% ===========================================================================
\section{Conclusion}
\label{sec:conclusion}

We have presented the Directive Grounding Hypothesis: that the intelligence
required to complete bounded software engineering tasks can be distributed
across model weights and context-provided skill definitions, with a
substitutability relationship between parametric capacity and skill
directiveness. This implies that aggressively quantized local models can
perform specialist agentic tasks when equipped with sufficiently directive
skill definitions—tasks that would fail with the same models under
minimal grounding.

The Agency framework is a practical instantiation of these principles,
providing 32 directive skill definitions organized as a complete virtual IT
organization. The framework's design—including Discovery Protocols, ordered
responsibility lists, named heuristic constraints, and explicit workflow
handoff rules—maps directly onto the contextual intelligence components
identified in our model.

The practical consequence is significant: teams can pursue capable AI-assisted
software development workflows using locally-hosted quantized models,
eliminating API dependencies and associated costs, provided they invest in
skill definition quality rather than raw model scale.

We invite the community to validate the Directive Grounding Hypothesis
through the experimental program outlined in Section~\ref{sec:future},
using the Agency skill library as a reference implementation.

% ===========================================================================
\section*{Acknowledgements}
The Agency framework is an open-source project. Contributions to skill
design, security validation (Basalt Shield), and CLI tooling are gratefully
acknowledged.

% ===========================================================================
\bibliographystyle{plainnat}
\bibliography{references}

\end{document}
